\documentclass[../../../include/open-logic-section]{subfiles}

\begin{document}

\olfileid{sth}{ordinals}{replacement}
\olsection{بدل رکھن دا اصول}

\olref[sth][ordinals][vn]{sec} وچ اسیں ترتیبی عدداں نوں متعارف کراؤن دا
جواز ایہ دتا سی کہ اسیں ایہناں نوں ترتیب دیاں قسماں، یعنی حُسنِ ترتیب
والیاں ساختاں دے معیاری نمائندے، منّ سکدے آں۔ ایس لئی سانوں ثابت کرنا
پئے گا کہ \emph{ہر حُسنِ ترتیب والی ساخت کسے ترتیبی عدد نال ہم ساخت اے}۔
ایس نال اسیں $\ordtype{A, <}$ نوں اوہ ترتیبی عدد $\alpha$ تعریف کر سکاں گے
جیہدے لئی $\tuple{A, <} \isomorphic \alpha$۔

بدقسمتی نال ہُن تک دتے مسلّمیاں نال ای اسیں مطلوبہ نتیجہ ثابت
\emph{نہیں کر سکدے}۔ (\olref[replacement][strength]{sec} وچ وجہ ویکھاں گے۔
ہُن لئی گل اینی اے کہ نہیں کر سکدے۔) سانوں اک نویں خیال دی لوڑ اے۔ اوہ ایہ اے:

\begin{axiom}[بدل رکھن دا خاکہ]
ہر فارمولے $\phi(x, y)$ لئی ہِیٹھلی گل اک مسلّمہ اے:
\begin{quote}
	کسے وی $A$ لئی، جے $(\forall x \in A)\lexists![y][\phi(x,y)]$، تاں
	$\Setabs{y}{(\exists x \in A)\phi(x,y)}$ موجود اے۔
\end{quote}
\end{axiom}
\noindent
وکھ کرن دے اصول وانگ ایہ اک خاکہ اے۔ ایہ لامتناہی گنتی وچ مسلّمے دیندا اے،
لامتناہی گنتی وچ وکھ وکھ $\phi$ وچوں ہر اک لئی اک۔ اینوں ایویں وی لکھ
سکدے آں (تے عام طور اُتے ایویں ای لکھدے آں):

\begin{defish}
ہر اوہ فارمولا $\phi(x,y)$ جیہدے وچ ’’$B$‘‘ نہیں آؤندا، اوہدے لئی
ہِیٹھلی گل اک مسلّمہ اے:
\[
\forall A[(\forall x \in A)\lexists![y][\phi(x,y)] \lif \exists B\forall y (y \in B \liff (\exists x \in A)\phi(x,y))]
\]
\end{defish}

پر پہلی واری ویکھن اُتے ایہ خاصا اُلجھیا فارمولا اے۔ بدل رکھن دے اصول
دا ہِیٹھلا سدھا نتیجہ شاید ساڈے بنیادی خیال نوں \emph{ہور صاف} طرح
بیان کردا اے:\footnote{ٹرم اوہ عبارت اے جیہڑی اکّو اک شے چُندی اے
(اوہدے آزاد متغیراں دیاں کوئی وی قدراں مقرر ہون اُتے)۔ مثلاً ’’$\emptyset$‘‘
اک ٹرم اے جیہڑی خالی سیٹ چُندی اے۔ ’’$\{x\}$‘‘ اک ٹرم اے جیہڑی
$x$ دا اک رکن والا سیٹ چُندی اے (بھانویں $x$ کوئی وی ہووے)۔
’’$x \cup y$‘‘ اک ٹرم اے جیہڑی $x$ تے $y$ دا اتحاد چُندی اے
(بھانویں اوہ کوئی وی ہون)۔}

\begin{cor}
کسے وی ٹرم $\tau(x)$ تے کسے وی سیٹ $A$ لئی ایہ سیٹ موجود اے:
\[
	\Setabs{\tau(x)}{x \in A} = \Setabs{y}{(\exists x \in A)y = \tau(x)}.
\]
\end{cor}

\begin{proof}
چونکہ $\tau$ اک \emph{ٹرم} اے، $\forall x \lexists![y][\tau(x) = y]$۔
ایس توں ہور وی صاف اے کہ $(\forall x \in A)\lexists![y][\tau(x) = y]$۔
سو $\Setabs{y}{(\exists x \in A)\tau(x) = y}$ بدل رکھن دے اصول نال موجود اے۔
\end{proof}
\noindent
ایس توں پتا لگدا اے کہ ’’بدل رکھنا‘‘ مسلّمے دا چنگا ناں اے۔ سیٹ $A$
دتا ہووے، تاں تسی نواں سیٹ $\Setabs{\tau(x)}{x \in A}$ بنا سکدے او،
$A$ دے ہر رکن دی تھاں اوہدا~$\tau$ تحت عکس رکھ کے۔ واقعی فنکشن تحت
سیٹ دے عکس والا نشان ورت کے اسیں $\funimage{\tau}{A}$ نوں
$\Setabs{\tau(x)}{x \in A}$ لئی لکھ سکدے آں۔

پر ضروری گل ایہ اے کہ $\tau$ اک \emph{ٹرم} اے۔ ایہ کوئی \emph{فنکشن}
(دا ناں) ہونا ضروری نہیں، جیویں \olref[sfr][fun][rel]{sec} وچ فنکشن
دا مطلب اے، یعنی مرتب جوڑیاں دا اک خاص سیٹ۔ آخر جے $f$ اوس معنی وچ
فنکشن اے، تاں سیٹ $\funimage{f}{A} = \Setabs{f(x)}{x \in A}$
بس $\ran{f}$ دا اک خاص ذیلی سیٹ اے، تے اوہدی موجودگی دی ضمانت
صرف~$\Zminus$ دے مسلّمے ای دیندے نیں۔\footnote{بس
$\Setabs{y \in \bigcup \bigcup f}{(\exists x \in A)y = f(x)}$ ویکھو۔}
اِس دے اُلٹ، بدل رکھن دا اصول ساڈے مسلّمیاں وچ اک \emph{طاقتور} وادھا اے،
جیویں اسیں \olref[sth][replacement][]{chap} وچ ویکھاں گے۔

\end{document}
